\chapter{Evaluation}

\section{Methodik}

\subsection{Studiendesign}
Zur Evaluation der entwickelten Anwendung \textit{petappoint} wurde 
eine quantitative Nutzerstudie durchgeführt. Als Messinstrument wurde 
die System Usability Scale (SUS) eingesetzt, die von John Brooke
entwickelt wurde \cite{brooke} und bis heute zu den am häufigsten eingesetzten 
Methoden zur Erfassung der wahrgenommenen Gebrauchstauglichkeit zählt 
\cite{sus}. Der SUS-Fragebogen umfasst zehn standardisierte Aussagen, die auf einer fünfstufigen Likert-Skala bewertet werden und als Ergebnis einen Score zwischen 0 und 100 liefern \cite{sus}.

Die Wahl des SUS begründet sich durch mehrere Faktoren: Erstens ist der Fragebogen technologieunabhängig und damit für mobile Anwendungen gleichermaßen geeignet wie für Webanwendungen. Zweitens ermöglicht 
der standardisierte Charakter des Instruments einen Vergleich der 
Ergebnisse mit bestehenden Benchmarks. Drittens ist der Fragebogen mit zehn Fragen bewusst kurz gehalten, was eine hohe Teilnahmebereitschaft begünstigt \cite{brooke}.

Ergänzt wurde der SUS-Fragebogen um einen weiteren Abschnitt zur plattformübergreifenden Wahrnehmung der Anwendung. Da \textit{petappoint} 
als plattformübergreifende mobile Anwendung konzipiert wurde und die Kompatibilität mit iOS und Android eine zentrale nicht-funktionale Anforderung darstellt (vgl. NFA\#01), zielt dieser Abschnitt darauf 
ab, die subjektiv wahrgenommene Qualität der Anwendung auf beiden Plattformen zu erfassen und gegenüberzustellen.

Die Studie richtet sich ausschließlich an Tierhalter, da diese die primäre Zielgruppe der Anwendung darstellen und den Buchungsprozess als Kernfunktion von \textit{petappoint} aktiv nutzen. Voraussetzung für die Teilnahme war, dass die Probanden die Anwendung zuvor 
eigenständig getestet hatten.

\subsection{Beschreibung der Datenerhebung}

Die Datenerhebung erfolgte mithilfe eines Online-Fragebogens, der über die Plattform Google Forms bereitgestellt wurde. Google Forms wurde aufgrund seiner einfachen Handhabung und der automatisierten Erfassung und Speicherung der Antworten gewählt.

HIER FEHLT DER PART WIE ICH DIE EVALUATION MACHE. ÜBER MOBILE WEB APP ODER ALS NATIVE MOBILE APP ÜBER 2 Smartphones.

Probanden bewusst gleichmäßig auf beide Plattformen aufgeteilt. 10 Personen testeten die Anwendung auf einem Android-Gerät und 
10 Personen auf einem iOS-Gerät. Im Anschluss an den Test wurde der Fragebogen den Probanden persönlich über einen QR-Code 
zugänglich gemacht, der auf den Google-Forms-Fragebogen verlinkte.

Die Teilnahme an der Studie war freiwillig und anonym. Es wurden keine personenbezogenen Daten wie Name oder E-Mail-Adresse erhoben. Die erfassten Daten wurden ausschließlich für wissenschaftliche 
Zwecke im Rahmen dieser Bachelorarbeit verwendet.

\subsection{Auswertungsverfahren}

Die Auswertung der erhobenen Daten erfolgt in drei Schritten: 
der Berechnung des SUS-Scores, dem Plattformvergleich zwischen 
iOS und Android sowie der qualitativen Auswertung der offenen 
Frage.

\subsubsection{SUS-Score}

Zur Berechnung des SUS-Scores wird die standardisierte Formel nach Brooke~\cite{brooke} angewendet. Dabei werden die Antworten der zehn SUS-Fragen wie folgt verarbeitet: Bei ungeraden Fragen (1, 3, 5, 7, 9) wird vom jeweiligen Antwortwert 
der Wert 1 subtrahiert. Bei geraden Fragen (2, 4, 6, 8, 10) wird der Antwortwert von 5 subtrahiert. Die so berechneten 
Teilwerte werden anschließend summiert und mit dem Faktor 2.5 multipliziert, woraus ein Gesamtscore zwischen 0 und 100 resultiert.

Die Interpretation des SUS-Scores erfolgt anhand der Bewertungsskala 
nach Sauro~\cite{sauro-interpretation}, der die Ergebnisse in 
Schulnoten einteilt:

\begin{itemize}
    \item 80--100: Note A (Sehr gut)
    \item 68--79: Note B (Gut)
    \item 51--67: Note C (Befriedigend)
    \item 26--50: Note D (Ausreichend)
    \item $\leq$ 25: Note F (Ungenügend)
\end{itemize}

Zur Auswertung wird zunächst der Durchschnitts-SUS-Score aller 20 Teilnehmer berechnet. Darüber hinaus werden die Scores der iOS- und Android-Nutzer getrennt ausgewertet und miteinander verglichen.

\subsubsection{Plattformvergleich}

Um die plattformübergreifende Qualität von \textit{petappoint} zu bewerten, werden die SUS-Scores sowie die Antworten auf die Fragen zur Plattformwahrnehmung (Fragen 13--15) getrennt nach Betriebssystem ausgewertet. Dabei werden die Mittelwerte der beiden Gruppen gegenübergestellt, um Unterschiede in der wahrgenommenen Benutzerfreundlichkeit zwischen iOS und Android zu identifizieren.

\subsubsection{Qualitative Auswertung}

Die Antworten auf die offene Frage werden thematisch gruppiert und auf wiederkehrende Muster untersucht. Ziel ist es, qualitative Erkenntnisse über Verbesserungspotenziale und Nutzerwünsche zu gewinnen, die durch die quantitativen SUS-Daten nicht erfasst werden.

\section{Ergebnisse}
\subsection{Anforderungsabdeckung der funktionalen und 
nicht-funktionalen Anforderungen}

Im Folgenden wird dargelegt, inwieweit die in Kapitel~2 Anforderungsermittlung
definierten funktionalen und nicht-funktionalen Anforderungen in der entwickelten Anwendung \textit{petappoint} umgesetzt wurden. Die Überprüfung erfolgt auf Basis der implementierten 
Funktionen der Anwendung.

\subsubsection{Funktionale Anforderungen}

\textbf{Tierhalter}

\renewcommand{\arraystretch}{1.5}
\begin{tabular}{|p{2cm}|p{3cm}|p{8cm}|}
\hline
\textbf{ID} & \textbf{Status} & \textbf{Anmerkung} \\
\hline
FA\#01 & Umgesetzt & Registrierung und Anmeldung über 
E-Mail und Passwort. \\
\hline
FA\#02 & Umgesetzt & Haustierprofile können angelegt, 
bearbeitet und gelöscht werden. \\
\hline
FA\#03 & Umgesetzt & Suche und Filterung nach 
Behandlungsart und Tierart möglich. \\
\hline
FA\#04 & Umgesetzt & Verfügbare Terminslots sind in einer 
Kalenderansicht einsehbar und buchbar. \\
\hline
FA\#05 & Umgesetzt & Gebuchte Termine sind einsehbar 
und können storniert werden. \\
\hline
\end{tabular}

\textbf{Tierarztpraxen}

\begin{tabular}{|p{2cm}|p{3cm}|p{8cm}|}
\hline
\textbf{ID} & \textbf{Status} & \textbf{Anmerkung} \\
\hline
FA\#07 & Nicht umgesetzt & Die Verwaltung des Praxisprofils 
sowie das Anlegen, Bearbeiten und Sperren von Terminslots 
wurden nicht implementiert. \\
\hline
FA\#08 & Nicht umgesetzt & Die Einsicht gebuchter Termine 
in einer Tages- und Wochenansicht wurde nicht 
implementiert. \\
\hline
\end{tabular}

\subsubsection{Zusammenfassung}

Die Analyse zeigt, dass alle funktionalen Anforderungen der Zielgruppe der Tierhalter vollständig umgesetzt wurden. Die funktionalen Anforderungen FA\#07 und FA\#08, die sich an Tierarztpraxen richten, wurden hingegen nicht implementiert und stellen ein Verbesserungspotenzial für die zukünftige Entwicklung dar. 

\subsection{SUS-Score und Interpretation}
\subsection{Plattformvergleich iOS vs. Android}
\subsection{Darstellungsfehler und qualitatives Feedback}

\section{Diskussion}
  \subsection{Interpretation der Ergebnisse}
  \subsection{Limitationen}